\chapter{Anhang A}
\label{chap:anhang_a}




\chapter{Nutzung generativer KI}
\label{chap:anhang_b}

<NAME TOOL\_1>; <BESCHREIBUNG ANLASS\_1>; <PROMPTS\_1>\\
<NAME TOOL\_2>; <BESCHREIBUNG ANLASS\_2>, <PROMPTS\_2>\\
…\\
<NAME TOOL\_n>; <BESCHREIBUNG ANLASS\_n>, <PROMPTS\_n>\\


\textcolor{gray}{
	Studierende müssen den Einsatz von generativer KI beim wissenschaftlichen Schreiben bei
	der Einreichung der Arbeit erklären. Die folgenden Hinweise beziehen sich nur auf den
	Schreibprozess und nicht auf die Verwendung von KI-Tools zur Analyse und Gewinnung von
	Erkenntnissen aus Daten als Teil des Forschungsprozesses:
	\begin{itemize}
		\item Generative KI und KI-gestützte Technologien sollten im Schreibprozess nur
		eingesetzt werden, um die Lesbarkeit und Sprache des Manuskripts zu verbessern.
		\item Die Technologie muss unter menschlicher Aufsicht und Kontrolle angewandt werden,
		und die Autoren sollten das Ergebnis sorgfältig überprüfen und bearbeiten, da KI
		autoritativ klingende Ergebnisse erzeugen kann, die falsch, unvollständig oder
		verzerrt sein können. Die Studierenden sind letztlich für den Inhalt ihrer Arbeit
		verantwortlich und rechenschaftspflichtig.
	\end{itemize}
	Text aus der Richtlinie der TU Braunschweig zur Nutzung generative KI
}

